\chapter{Hipótesis y Objetivos del Proyecto}

\section{Hipótesis}
Se plantea que es posible diseñar e implementar un sistema de control ESD de bajo costo que permita medir 
la resistencia cuerpo--tierra en el rango de $500~\mathrm{k}\Omega$ a $40~\mathrm{M}\Omega$, con una precisión de $\pm10\%$, 
integrando funcionalidades de identificación, monitoreo y registro, de manera que el sistema sea funcionalmente 
equivalente a soluciones comerciales en términos de verificación ESD para entornos de laboratorio.

\section{Objetivos}

        \subsection{Objetivo General}
            Diseñar e implementar un sistema inteligente de control ESD de bajo costo, 
            que integre múltiples funcionalidades orientadas a entornos industriales y de laboratorio.

        \subsection{Objetivos Específicos}
                \begin{itemize}
                    \item Diseñar e implementar un sistema de medición de resistencia cuerpo--tierra en el rango de 
                    $500~\mathrm{k}\Omega$ a $40~\mathrm{M}\Omega$, con una precisión aproximada de $\pm 10\%$.

                \item Desarrollar la arquitectura hardware del sistema, incluyendo la PCB, fuente de alimentación 
                y técnicas de medición para alta impedancia, considerando la reducción de corrientes de fuga.

                \item Integrar funcionalidades de monitoreo e interacción, incluyendo sensores ambientales, 
                detección de usuario, identificación mediante RFID e interfaz gráfica para visualización de resultados.

                \item Implementar un sistema de registro y generación de reportes, y validar el funcionamiento del 
                sistema en el laboratorio TestChip de Intel Costa Rica.
            \end{itemize}

\section{Alcances y Limitaciones del Proyecto}

